\section{How to think about evaluation}

\subsection{What's the point of evaluation?}

老师最后回到本讲开头：评价没有唯一真理，而是服务不同问题。用户或公司需要为具体 use case 做采购选择；研究者想测 raw capability；商业与政策评估关心 benefits 和 harms；模型开发者还需要失败反馈来改进系统。目的不同，数据分布、metric、judge、工具权限与成本约束就不应相同。

\begin{knowledgebox}{老师强调：没有 one true evaluation}
课程给出的四类目的不是要选出唯一正确答案，而是要求每份报告先说明自己在回答什么问题。购买决策、科学测量、风险评估和开发反馈可以共享部分 benchmark，却不能默认共享同一 ranking；分数只有在目标和规则匹配时才有意义。
\end{knowledgebox}

\subsection{What are we evaluating?}

Pre-foundation model 时代我们经常评价 methods：固定数据、固定 split、比较算法。今天我们更多评价 models/systems：任何训练数据、prompting、tool scaffold 都可能进入系统，规则必须明说。NanoGPT speedrun 是少数仍明确固定数据与目标 loss、比较达到目标所需 compute time 的 method evaluation。

\begin{figure}[H]
\centering
\includegraphics[width=0.74\textwidth]{karpathy-nanogpt-speedrun.png}
\caption{nanoGPT speedrun：少数仍在评价 method 的例子，固定数据和目标 validation loss，比较 compute time。}
\figsource{课程源引用图，本地化后供 CS336 Lecture 12 使用。}
\end{figure}

\begin{knowledgebox}{读图：method eval 和 system eval 的区别}
nanoGPT speedrun 固定数据和目标 validation loss，比较达到目标所需时间，因此鼓励算法和工程优化。Arena 或 SWE-Bench 则更像 system eval：用户关心哪个系统好用，不一定关心改进来自模型、数据、prompt 还是 scaffold。
\end{knowledgebox}

这一区分解释了许多论文和产品报告为什么看起来互相矛盾。研究论文想证明一个方法更好，就应该固定数据、计算预算、评测规则，尽量隔离变量；产品评测想告诉用户哪个系统好用，就可以允许工具、检索、prompting、post-training 和产品策略一起发挥作用。二者都合理，但不能混在同一张表里解释。

\begin{knowledgebox}{术语消化：methods vs models vs agents}
\begin{tabular}{L{0.2\textwidth}L{0.36\textwidth}L{0.34\textwidth}}
\toprule
对象 & 评价什么 & 例子 \\
\midrule
Methods & 算法或训练方法，在固定数据和规则下比较。 & nanoGPT speedrun。 \\
Models & 训练好的模型能力和偏好表现。 & MMLU、Arena、HLE。 \\
Agents/systems & 模型加 scaffold、工具、记忆和执行循环。 & SWE-Bench、TerminalBench、CyBench。 \\
\bottomrule
\end{tabular}
\end{knowledgebox}

\begin{importantbox}{讲义提醒：评价规则本身就是研究对象}
当模型越来越强，旧 benchmark 会饱和，评测者会加难；当模型会调用工具，评测者要决定是否允许工具；当模型会长链推理，评测者要决定是否限制时间和 token；当模型会影响真实用户，评测者要考虑隐私和生态有效性。Evaluation 不是训练结束后的打分脚本，而是 AI 系统发展的一部分。
\end{importantbox}

